\section{Introduction}\label{sec:introduction}

For a univalent function, the derivative can become large or small near
the boundary even when the image domain is bounded. Integral means
measure the rate of this concentration. Their universal spectrum records
the largest possible growth exponent over a class of conformal maps;
it brings questions about derivatives, coefficients, and the geometry
of image boundaries into a common framework.

Let \(\DD=\{z\in\C:|z|<1\}\), and let \(\SSS\) consist of all
injective holomorphic maps \(f:\DD\to\C\) with \(f(0)=0\) and
\(f'(0)=1\). Write \(\SSS_b\) for the subclass of maps with bounded
image. The bound on the image may depend on \(f\), and no boundary
regularity is imposed. For \(p\in\R\) and \(0\le r<1\), define
\begin{align*}
 M_p[f'](r)&=\frac1{2\pi}\int_{-\pi}^{\pi}|f'(re^{i\theta})|^p\,d\theta,\\
 \beta_f(p)&=\limsup_{r\uparrow1}
 \frac{\log M_p[f'](r)}{\log(1/(1-r))},
 &B_b(p)&=\sup_{f\in\SSS_b}\beta_f(p).
\end{align*}
The derivative of a univalent map never vanishes, so all real powers
are defined. The supremum in \(B_b\) is taken after each map's limiting
exponent. Multiplying a map by a nonzero constant leaves that exponent
unchanged, so normalization of the derivative does not change the
bounded-class spectrum.

\subsection{The conjecture and the uniform bound}

Kraetzer's experiments with conformal maps onto bounded domains led
to the prediction
\begin{equation}\label{eq:kraetzer}
 B_b(p)=
 \begin{cases}
 p^2/4,&|p|\le2,\\
 |p|-1,&|p|\ge2.
 \end{cases}
\end{equation}
The quadratic part is stated in his 1996 paper and explicitly
restated in his thesis \cite{Kraetzer,KraetzerThesis}; the full
quadratic--linear formulation also appears in Beliaev's discussion
of numerical spectra \cite{Beliaev}. We prove a strict upper bound
at the negative interior exponent \(p=-1\).

\begin{theorem}[Uniform inverse-first-power bound]\label{thm:main}
There exist constants \(0<\varepsilon<1/4\) and \(C<\infty\) such that
\begin{equation}\label{eq:uniform}
 M_{-1}[f'](r)\le C(1-r)^{-1/4+\varepsilon}
 \qquad\text{for every }f\in\SSS\text{ and every }1/2\le r<1.
\end{equation}
Consequently \(B_b(-1)<1/4\), and Kraetzer's conjectured identity
\eqref{eq:kraetzer} is false.
\end{theorem}

Both constants are independent of the map, including its image size.
The proof gives existence of a positive gap but does not evaluate it
numerically. The conclusion concerns the inverse first power; it does
not settle the proposed value at \(p=1\).

\subsection{Context and the role of strictness}

The distinction between bounded and unrestricted disk maps matters.
Carleson and Jones connected coefficient problems with conformal
dimension \cite{CJ}. Makarov's comparison theorem identifies the
unrestricted spectrum with the maximum of the bounded spectrum and
the Koebe contribution \(3p-1\) \cite[Theorem 5.4]{Makarov}.
In particular, the bounded spectrum in \eqref{eq:kraetzer} should
not be replaced by the unrestricted disk spectrum at positive exponents.
Hedenmalm and Shimorin developed a weighted Bergman-space approach
to bounds for these spectra, keeping this distinction explicit
\cite{HedenmalmShimorin}. Their upper bound at \(p=-1\) for the
unrestricted normalized disk class is \(0.403\)
\cite[Table 1]{HedenmalmShimorin}. Sola subsequently obtained the
bound \(0.388\) by related Bergman-space methods
\cite{Sola}; see also the author's account in
\cite[p.~6]{SolaThesis}. These estimates remain above the
conjectured value \(1/4\); Theorem~\ref{thm:main} establishes a
uniform exponent strictly below it.

Random conformal snowflakes provide another approach. Beliaev and
Smirnov constructed random conformal maps whose averaged spectrum
approaches the universal spectrum at any prescribed exponent, and
related lower bounds to positive integral operators
\cite[Sections 3--4]{BeliaevSmirnov}. Beliaev further investigated
these operators numerically, distinguishing the numerical estimates
from rigorous error control \cite{Beliaev}. The common feature with
the present argument is the use of a law on conformal maps and a
positive operator. Here the operator samples inverse derivatives
through affine changes of base point, and its growth controls an
upper bound uniformly over the entire class \(\SSS\).

The affine-law and saddle-to-maximum construction is also used for
inverse-square means in the companion article
\emph{Brennan's conjecture and sharp inverse-square integral means}
\cite{BrennanMethod}. The two arguments share the compact
normalization of half-plane maps, dyadic sampling, and a comparison
of critical points over a common target. We prove every
inverse-first-power estimate and every transport assertion needed
here; the companion's inverse-square endpoint theorem is not an
input.

The obstacle specific to the present theorem is strictness at the
candidate exponent \(1/4\). A non-strict comparison of saddle and
maximum weights only rules out larger exponents. We retain the exact
operator growth parameter when constructing the probability law.
For almost every target, distinct critical points have distinct
values, so the rank pairing sends each saddle to a strictly higher
maximum. This strict increase forces a vanishing Jacobian in the
remaining equality case. A harmonic lower barrier then excludes
that case. The strict increase can be arbitrarily small from one
pair to another; no common pointwise gain is required.

\subsection{Proof outline}

Section~\ref{sec:operator} transfers disk maps to the compact family
of univalent maps on the upper half-plane \(\HH=\{x+iy:y>0\}\),
normalized by \(F(i)=0\),
\(F'(i)=1\). A positive operator \(L\) samples \(1/|F'|\) at two
affine subscales. The exponential growth rate of \(\|L^n\|\), denoted
by \(2^\beta\), bounds all the disk means in Theorem~\ref{thm:main}.
If \(\beta\ge1/4\), Section~\ref{sec:law} constructs a probability
law whose inverse-derivative weighted rebasings scale exactly as
\(y^{-\beta}\). Differentiating this identity fixes two moments of
the logarithmic derivative.

For a target \(\xi\in\C\), the potential
\(y^{(\beta+2)/3}/|F(x+iy)-\xi|\) turns those moment identities into
a geometric comparison. Section~\ref{sec:escape} proves compactness
of its positive superlevels for almost every map and target.
Section~\ref{sec:pairing} identifies its critical points through a
smooth map \(G_F\), proves the target-generic absence of degeneracies
and ties, and constructs the measurable rank pairing. The nullity of singular
values is the planar instance of Sard's theorem \cite{Sard}; injectivity
supplies the separate exclusion of tied critical values. The planar
index count follows the tangent-turning argument of Hopf \cite{Hopf}.
The moment identities make the expected normalized signed Jacobian
of \(G_F\) positive when \(\beta>1/4\), and zero when \(\beta=1/4\).

Section~\ref{sec:transport} changes variables through \(G_F\) and
compares the paired weights on expanding source rectangles. It also
transfers exceptional-target nullity to the fixed normalized point,
where area-genericity alone would be insufficient.
The comparison forces the expected signed Jacobian to be nonpositive;
strictness of the pairing forces the Jacobian itself to vanish almost
surely in the remaining equality case.
Section~\ref{sec:equality} rules out the vanishing-Jacobian equality
case and converts \(\beta<1/4\) into the uniform estimate.
